\sheet{1}{Architecture and memory effects}{1h30}
\section{Preparation}
Define \textbf{register}, \textbf{instruction dependency}, \textbf{latency},
\textbf{throughput}, \textbf{cache line}, \textbf{spatial locality} and
\textbf{temporal locality}. Does a 3~GHz clock imply three billion completed
instructions per second? Identify the missing information.

The source files are in \code{Exercices/01-architecture}. Each file contains
one program: the functions to complete and a supplied \code{main} that checks
small inputs before measuring execution time. Edit the functions marked TODO.
An unfinished function reports its name and stops; another exercise is not needed
to run the current one.

Run commands from the \code{PW} directory. Configure once:
\begin{lstlisting}[style=shell]
cmake -S . -B build -DCMAKE_BUILD_TYPE=Release
\end{lstlisting}
Rebuild the relevant target after every source change. With no MPI or OpenMP
installation, add \code{-DPW_ENABLE_MPI=OFF -DPW_ENABLE_OPENMP=OFF} to configuration.
The numerical arguments and the expected checks are explained in each section.

\section{Experiments}
\exercise{A1}{One accumulator and four independent accumulators}
Edit \code{A1-sum.cpp}. The supplied \code{sum_one} adds an array using one
accumulator. Each addition needs the preceding addition's result. Four separate
accumulators can expose independent additions to the same core; this still uses
one software thread.
\begin{enumerate}
\item Draw four successive updates of one accumulator. Assume addition latency
is 4 cycles and the processor can start one independent addition per cycle.
Derive the sustained addition rates for one chain and four chains at 3~GHz.
Ignore loads, loop control, startup and final combination in this model.
\item Implement \code{sum_four}. Initialise four sums, process four elements
per main-loop iteration, then handle the remaining elements. Combine the sums
and return the result. Trace lengths 0, 1, 3, 4 and 5 to verify coverage and bounds.
\item Build and run:
\begin{lstlisting}[style=shell]
cmake --build build --target a1_sum -j
./build/bin/a1_sum 1048576 7
\end{lstlisting}
The arguments are array length and trial count. The program checks empty arrays,
short arrays and tails, then prints individual times, medians and ranges.
Checks and printing are outside the measured kernel intervals.
\item Compare \code{one} and \code{four} for a small array and an array larger
than the last-level cache but comfortably within RAM. Record both medians and
$T_{\mathrm{one}}/T_{\mathrm{four}}$. Explain why limited data delivery may hide
an improvement in arithmetic throughput. Generated machine instructions may also
differ from the source-level model.
\item The program also evaluates \code{{1e16,1.0,-1e16,1.0}}. Its exact sum is 2.
Trace the rounding for both functions. Explain why different results can occur
without threads or data races. The ordinary input tests use exactly representable
values; this separate example illustrates sensitivity to grouping.
\end{enumerate}

\exercise{A2}{Reading a matrix by rows and columns}
Edit \code{A2-matrix.cpp}. A square matrix is stored in one contiguous array:
row $r$, column $c$ is at position $r\,\mathrm{side}+c$. The supplied
\code{matrix_rows} visits consecutive elements. A column traversal reads the
same values in another order, with a jump of \code{side} elements between rows.
\begin{enumerate}
\item Draw a $4\times4$ matrix and label its storage offsets 0 to 15. List the
first eight offsets for each traversal. Assume 8-byte doubles and aligned
64-byte cache lines: mark the values brought in by the first line fetch.
\item Implement \code{matrix_columns} with the column loop outside the row
loop. Do not copy or transpose the data. Return the sum of all entries.
\item Build and run:
\begin{lstlisting}[style=shell]
cmake --build build --target a2_matrix -j
./build/bin/a2_matrix 512 7
\end{lstlisting}
The arguments are matrix side and trials. The checks include sides 0, 1, 3, 17
and 65. Repeat the timings for sides 64, 512 and 2048, reducing sizes if RAM is
limited. Compare row and column medians for identical stored data.
\item Explain the results using immediate use of neighbours and the interval
between later uses of the same cache line. A small cached matrix need not show
the same difference as a large one. Check the actual cache-line size before
applying the 64-byte model to this machine.
\end{enumerate}

\exercise{A3}{Stride and cost per useful access}
Edit \code{A3-stride.cpp}. Stride 1 reads every element; stride 3 reads indices
0, 3, 6, and so on. Increasing stride changes both the access pattern and the
number of additions, so total time alone is not a fair measure of cost per access.
\begin{enumerate}
\item Implement \code{strided_sum} for a positive stride. Return zero for empty
input and stop before index $n$. For $n=10$, list the visited indices for strides
1, 3 and 16. Derive the visit count $V$; handle $n=0$ separately.
\item Build and run:
\begin{lstlisting}[style=shell]
cmake --build build --target a3_stride -j
./build/bin/a3_stride 1048576 7
\end{lstlisting}
The arguments are array length and trials. Small-input checks include empty
input and a stride larger than the array. The timing run uses strides
1, 2, 4, 8, 16, 32 and 64 at one fixed length.
\item For median time $t$ seconds, calculate $10^9t/V$ nanoseconds per visit
and $8V/(10^9t)$ useful GB/s for 8-byte doubles. Compare your calculation with
the output. Explain why useful bytes do not count all physical cache-line or
DRAM traffic.
\item Plot time per visit against stride. Explain possible effects of cache
capacity, prefetching and unused values on fetched lines. For 17 doubles and
stride 4, predict the visited indices, number of visits and useful bytes.
\end{enumerate}

\exercise{A4}{Optional: tiled transpose}
Edit \code{A4-transpose.cpp}. A transpose writes $A_{r,c}$ to $B_{c,r}$ in a
separate output array. A simple traversal can read consecutive input values while
writing distant output values. Small tiles can retain useful lines from both arrays.
\begin{enumerate}
\item Read the supplied \code{transpose_reference}. Implement
\code{transpose_tiled} using loops over tile origins and loops inside a tile.
Clamp the final tile to the matrix boundary; never read beyond the input.
\item Build and run:
\begin{lstlisting}[style=shell]
cmake --build build --target a4_transpose -j
./build/bin/a4_transpose 1024 7
\end{lstlisting}
The arguments are side and trials. The program compares every output element,
including partial tiles, before timing the simple traversal and tile sides
8, 16, 32 and 64.
\item Explain why a tile that is too large or too small can lose the expected
benefit. Compare identical matrix sizes and inspect variation before choosing
one tile size from close measurements.
\end{enumerate}
